\section{Introduction}
\subsection{Empathy}
The Tragedy of the Commons, introduced by \cite{hardinTragedyCommonsPopulation1968}, describes how population growth and resource demand can lead to overexploitation of shared resources. Amid the current environmental crisis, understanding psycho-social properties and organisational frameworks to limit the risk of the Tragedy in humans is becoming existentially crucial. This research seeks to identify theoretical implications of the current understanding of empathy on resource consumption and its fairness with the use of a decentralized model-free multi-agent reinforcement learning framework.
Economists study this issue through the paradigm of common-pool resource dilemmas, where groups decide how much of a shared resource to extract, but individuals can exploit more or less of the resource than the agreed-upon amount \cite{patelModellingCooperationCompetition2021}. Common-pool resources are defined by two properties \cite{ostromRevisitingCommonsLocal1999,perolatMultiagentReinforcementLearning2017}:
(1) Non-exclusivity: agents cannot easily be excluded from consuming the resources.
(2) Subtractability: agent's consumption reduces the availability to others.
Classic economic theory predicts that rational agents, aiming to maximise their marginal utility, converge towards a Nash Equilibrium \cite{nashEquilibriumPointsPerson1950}. Cumulative exploitation leads to resource depletion and suboptimal outcomes for individuals and the group \cite{patelModellingCooperationCompetition2021}. This reflects crucial tensions between self-interest and collective welfare, and between short-term and long-term interests \cite{vanlangeHowWeCan2008}.


\subsection{Multi-agent reinforcement learning}
Multi-agent reinforcement learning (MARL) has increasingly been used as a computational framework for modelling socio-cognitive phenomena in interactive environments \cite{gronauerMultiagentDeepReinforcement2022} and has even been applied to model multi-agent behaviors in a tragedy of the commons scenario \cite{perolatMultiagentReinforcementLearning2017}.
 To study common-pool resources and other social dilemmas, decentralized learning Model-Free (MF) agents are frequently used, thus framing these dilemmas as reward maximization problems \cite{perolatMultiagentReinforcementLearning2017,wangEvolvingIntrinsicMotivations2018,zhuDeepMultiagentReinforcement2019}.
 MF agents incrementally update their policies based on direct reward experiences, linking them to procedural learning \cite{kosterModelfreeConventionsMultiagent2020,massiModelbasedModelfreeReplay2022}. In contrast, model-based agents integrate an internal representation of the environment, allowing them to simulate outcomes and adapt strategies more flexibly to changes. Although this approach allows simulation of decision-making and planning capabilities, it demands greater computational resources and assumes advanced agent capabilities \cite{massiModelbasedModelfreeReplay2022}.

MF-MARL simulations have recently been used to understand the dilemmas of common-pool resources. \cite{perolatMultiagentReinforcementLearning2017} demonstrated that allowing agents to stop others for a set duration, thus temporarily limiting the number of active foraging ones, allowed the system to reach an equilibrium. They observed three phases of the system: naïvety (i.e., increased utility but decreased sustainability), tragedy (i.e., declines in sustainability, peace, and utility), and maturity (i.e., utility and sustainability improved steadily despite reduced freezing action selection and stable equality). Similarly, \cite{zhuDeepMultiagentReinforcement2019} experimentally demonstrated that larger groups depleted resources faster, but that higher valuation of sustainability parameter in the reward function could mitigate that effect. Finally, \cite{kosterModelfreeConventionsMultiagent2020} observed that agents with heterogeneous preferences formed bottom-up conventions. They prioritized maximizing others' optimality over their hypothetical gain. This showcases emergent cooperative behaviors. To our knowledge, no existing work has yet proposed a way to model empathetic processes within a MARL framework. We now turn to examining key properties that may influence the likelihood of a Tragedy of the Commons, with a particular focus on mechanisms like empathy, through the observation of others' states and the valuation of their utility.

\subsection{Determinants of the Tragedy}
\subsubsection{Top-down organizational determinants.}
From a governance system and economic perspective, effective governance requires balancing resource access with sustainable consumption incentives, addressing both collective needs and individual use \cite{ostromRevisitingCommonsLocal1999,ostromGeneralFrameworkAnalyzing2009}. Ostrom \cite{ostromGeneralFrameworkAnalyzing2009} identified five interacting conditions of effective management of common-pool resources at the local and regional levels. First, users should benefit from organizing collectively. Second, reliable and valid indicators of resource conditions should enable users to assess long-term costs and benefits. Third, users must anticipate that the collective benefits outweigh the personal costs. Fourth, there should be sufficient trust among users or low-cost mechanisms to monitor and sanction non-compliance. Finally, users should actively participate in the decision-making process. In addition, \cite{ostromGeneralFrameworkAnalyzing2009} introduced a comprehensive framework to analyze the sustainability of socioecological systems. This framework links users, governance systems, resource systems, and resource units with their outcomes and interactions while considering the ecological, social, and economic contexts.

\subsubsection{Group and interaction level determinants}
Although institutional frameworks are essential for managing individual resource consumption, social interactions at the local level also play a critical role in promoting sustainable resource use. In their survey of Maasai pastoralists, \cite{rabinovichProtectingCommonsPredictors2020} observed that group membership was associated with the prioritization of collective outcomes and resource protection. Similarly, \cite{cuadradoRoleProsocialnessTrust2017} observed that in a multiplayer farming game, those who could communicate and had a secondary common goal of maximizing global profit chose less selfish irrigation strategies and obtained greater individual gains. Furthermore, \cite{decaroSynergisticEffectsVoting2015,decaroMotivationalFoundationsCommunication2021} observed that in a blind common-pool resources study, democratic decisions and social rewards favored the respect for negotiated rules and performance. They increased needs-satisfaction, intrinsic motivation, and trust with effects lasting even after communication was interrupted. The humanistic rational choice theory explains these results by proposing that a virtuous cycle of conditional cooperative self-governance can emerge from the satisfaction of three core psychosocial processes: (1) fundamental need satisfaction, (2) institutional internalisation and acceptance, and (3) group cohesion \cite{decaroMotivationalFoundationsCommunication2021}. Along this view, fundamental needs are evolved socio-psychological elements essential to human well-being: procedural justice (fairness), self-determination, security, efficacy and understanding, social acceptance and meaningful social relationships, and fair distribution of costs/benefits of governance.

\subsubsection{Individual-level determinants}
In the context of behavioral interventions, researchers have identified various psychosocial variables—such as values, beliefs, and social norms—as key predictors of pro-environmental behavior \cite{fornaraExtendedValueBeliefNormTheory2020,guoResidentialElectricityConsumption2018,liWhatInfluencesIndividuals2019,sussmanProEnvironmentalValuesMatter2016}. To integrate these determinants, Thiermann and Sheate \cite{thiermannMotivatingIndividualsSocial2020} proposed the two-pathway model of pro-environmental behavior that connects the Value-Belief-Norm (VBN) model \cite{sternValuebeliefnormTheorySupport1999}, the Theory of Planned Behavior (TPB; \cite{ajzenTheoryPlannedBehavior1991}), habits, and empathetic processes. The normative pathway links values and intentions to behavior. In contrast, the relational pathway highlights the influence of empathy, compassion, and connectedness with nature.
In a related perspective, \cite{moreno-casasEthicsCareTragedy2021} suggest that the ethics of care, grounded in empathic recognition of others' needs, can mitigate the tragedy of the commons by fostering cooperative behavior and concern for collective well-being. Similarly, the dual interest theory developed by \cite{lynneTheoreticalFoundationEmpathy2016} posits that individuals may derive utility through empathic processes, allowing empathy to curb self-interest and promote sustainable behavior.
From an evolutionary standpoint, empathy may have been selected as an adaptive trait that enhanced group-based foraging strategies and cooperation, thereby supporting human survival and reproduction (\cite{tomaselloTwoKeySteps2012}). Over time, environmental pressures favored empathy and affective signaling as mechanisms for promoting group cohesion, moral concern, and the sharing of resources \cite{cuffEmpathyReviewConcept2016,decetyEmergenceEmpathyDevelopmental2021,tomaselloOriginsHumanCooperation2013}. Having established empathy as a potentially key determinant of pro-environmental behaviors, the next section presents a theoretical framework to define it.



\subsection{Defining empathy}
The literature review above emphasized how empathy can be considered as an adaptive trait that may have enhanced group-based cooperation. Nevertheless, in the psychology literature, multiple definitions of empathy exist, and it is important to be clear about what we mean by empathy here before being able to formalize a tentative mechanism for its influence on reinforcement learning processes.

\subsubsection{The object of empathy: what are emotions?}
Emotions, the primary objects of empathetic processes, can be defined as categorical entities (e.g., fear, joy) composed of five components: phenomenological experience, physiological response, expressive behaviors, behavioral tendencies, and emotivational goals \cite{rosemanEmotionalBehaviorsEmotivational2011}. Alternatively, emotions can be described dimensionally, along axes such as valence and activation within a circumplex model \cite{russellCoreAffectPsychological2003,rosemanEmotionalBehaviorsEmotivational2011}. Appraisal theories propose that emotions result from rapid evaluations of stimuli based on criteria such as goal relevance and perceived control \cite{broschCommentAppraisingBrain2013,moorsPowerGoalDirectedProcesses2017}. Thus, depending on agents' goals, both positive and negative emotions may be induced by resource consumption and depletion. Emotional expressions, especially those associated with valence (e.g., distress or satisfaction), can act as social signals and influence others' decisions regarding resource sharing, particularly in common-pool dilemmas \cite{lynneTheoreticalFoundationEmpathy2016}.

\subsubsection{Processing others' emotions: what is empathy?}
Empathy is a complex construct that lacks a universally accepted definition and is often conflated with related concepts such as sympathy, emotional contagion, compassion, and prosociality \cite{cuffEmpathyReviewConcept2016}. Nonetheless, it is generally agreed to involve both understanding and sharing the emotional state of another person, with the recognition that the other is the source. Empathy arises from the interaction between trait capacities and situational influences, combining automatic emotion-eliciting processes and top-down control \cite{cuffEmpathyReviewConcept2016}. This results in an emotional experience aligned with the perception of another person's emotion, while maintaining self-other distinction.

\subsubsection{Three components of empathy: cognitive, affective, and motivational}
Empathy is commonly described as comprising three components, each supported by distinct neural systems:
\begin{itemize}
	\item \textbf{Cognitive empathy}: The capacity to understand or imagine what another person is feeling.
	\item \textbf{Affective empathy}: The capacity to share or resonate with the emotional state of another person \cite{bosnjakovicEmpathyConceptsTheories2018,dewaalMammalianEmpathyBehavioural2017}. Affective empathy is distinguished from emotional contagion by the awareness that the emotion originates in another person, not the self. The role of mirror neurons in empathy remains debated, with mixed empirical results \cite{bekkaliPutativeMirrorNeuron2021,boniniMirrorNeurons302022}. Both cognitive and affective empathy induce similar emotion-specific neural activations, with differences related to self-other differentiation \cite{bernhardtNeuralBasisEmpathy2012,dvashTheoryMindEmpathy2014}.
	\item \textbf{Motivational empathy (empathic concern)}: The drive to care for others' well-being, often considered a third and independent component, or as a form of prosocial motivation encompassing sympathy and compassion \cite{dewaalMammalianEmpathyBehavioural2017,decetyEmergenceEmpathyDevelopmental2021}.
\end{itemize}

While cognitive and affective empathy allow us to have a correct estimate of the other's utility, they do not intrinsically influence our behavior. Following utility based calculus theory of action selection and the empathic-altruistic hypothesis \cite{Batson01041991,batson2022empathy}, the empathic concern or motivational component is the key to drive behavior from the proper estimation of the others' states.

\subsection{Theoretical Hypotheses}
Drawing from social psychology models, economic theories \cite{lynneTheoreticalFoundationEmpathy2016,thiermannMotivatingIndividualsSocial2020}, and findings from MARL research on prosocial behavior \cite{yuEmotionalMultiagentReinforcement2015}, four theoretical hypotheses (TH) are proposed, organized around two objectives.

Objective 1: To study the independent and combined contributions of the components of empathy. We expect that both cognitive empathy and motivational empathy will independently reduce resource overconsumption, thereby delaying depletion and extending episode length, with their combination producing the strongest effect on sustainability (TH1). Additionally, we expect that both components will independently promote fairer resource access, with their interaction yielding the most equitable distribution outcomes (TH2).

Objective 2: To examine the parametric modulation of the motivational empathy component across a continuous range of altruistic valuation. We expect a positive monotonic relationship between the degree of altruistic motivation and sustainability, such that higher levels of motivational empathy will be associated with longer episodes before resource depletion (TH3). Similarly, we expect that increasing altruistic motivation will produce progressively more equitable patterns of resource consumption across agents (TH4).

